% ============================================================
%  Chapter 6: Discussion
%  File: chapters/06_discussion.tex
%  Bibliography: chapters/refs_06_discussion.bib
% ============================================================

\chapter{Discussion}
\label{ch:discussion}

The results in Chapter~\ref{ch:results} raise several questions about
the nature of the evolved structures, the boundary of the architecture's
capabilities, and the relationship between \modnet{} and existing
approaches to neuroevolution. This chapter interprets the main findings,
discusses limitations and threats to validity, and proposes directions
for future work.

% ------------------------------------------------------------
\section{Interpretation of the Main Findings}
\label{sec:disc-interpretation}
% ------------------------------------------------------------

\subsection{Layer-Free Computation Is Feasible for Digital Logic}
\label{sec:disc-layer-free}
% ------------------------------------------------------------

The most striking result of this work is that a fully layer-free
architecture, with no predefined modules and no separate output layer,
solves \textbf{eleven of eleven} Boolean and digital-circuit tasks to
exact $100\%$ accuracy (Table~\ref{tab:main-results}). This includes
tasks that are notoriously difficult for neural networks: the full
adder, the 2-bit adder with three output bits, the 4-bit comparator
with one-hot encoding, the 4-to-1 multiplexer with six inputs, and the
4-bit sorting network.

This is not obvious a priori. The layered architecture of modern neural
networks is often justified by the fact that it makes the network's
computation ``deep,'' allowing the composition of many simple
non-linearities into a complex function. A layer-free architecture has
no notion of depth. Yet our results show that depth-like computation
emerges naturally from the recurrent flow of information through a
cyclic graph.

A useful analogy is the difference between a feedforward circuit and a
general program with loops. A feedforward circuit computes a function
of bounded depth; a program with loops can compute functions of
unbounded depth by iterating. \modnet{} is closer to the latter: each
time step is like one iteration of a loop, and the network's output is
a function of the input after $T$ iterations. This is not a new
observation---it is the basis of recurrent neural networks
\citep{siegelmann1995computational}---but our results show that it can
be achieved without any predefined layer structure and without any
modular scaffold.

\subsection{Recurrence Emerges and Scales with Task Complexity}
\label{sec:disc-recurrence}
% ------------------------------------------------------------

Every successful network contained recurrent self-loops. Crucially, the
number of self-loops scales with the temporal complexity of the task,
as shown in Table~\ref{tab:recurrence-disc}.

\begin{table}[h]
  \centering
  \small
  \caption{Recurrence count scales with task complexity. XOR-2, a purely
    combinatorial problem, requires a single self-loop. Gaussian and
    multiplication, both continuous, require ten to fourteen.}
  \label{tab:recurrence-disc}
  \begin{tabular}{lccc}
    \toprule
    \textbf{Task} & \textbf{Type} & \textbf{Self-loops} & \textbf{Nodes} \\
    \midrule
    XOR-2            & Boolean    &  1 &  8 \\
    XOR-3            & Boolean    &  6 &  9 \\
    Parity-4         & Boolean    &  5 & 10 \\
    Parity-6         & Boolean    &  5 &  8 \\
    Full Adder       & Boolean    &  5 & 26 \\
    2-bit Adder      & Boolean    &  6 & 26 \\
    Comparator       & Boolean    &  6 & 27 \\
    MUX 4-to-1       & Boolean    &  4 & 26 \\
    Sort-4           & Boolean    &  6 & 26 \\
    \midrule
    $\sin(2\pi x)$   & Regression & 12 & 18 \\
    $\sin(4\pi x)$   & Regression & 11 & 13 \\
    $x^2 + y^2$      & Regression & 12 & 26 \\
    Gaussian         & Regression & 14 & 27 \\
    Multiplication   & Regression & 10 & 13 \\
    \bottomrule
  \end{tabular}
\end{table}

The pattern is clear. Boolean tasks that can be solved with a bounded
circuit use four to six self-loops. Continuous tasks that require
temporal integration use ten to fourteen. This suggests that recurrence
is a \emph{structural necessity} for tasks that require temporal
integration, and that \modnet{} discovers this necessity without any
architectural prior.

No mechanism in our algorithm explicitly favors recurrence. Self-loops
emerge because the unrestricted topology allows them, and evolutionary
pressure discovers them to be useful. If recurrence were useless, it
would have been pruned. Instead, every successful network contains it.

This is consistent with the biological observation that recurrent
connectivity is ubiquitous in cortical circuits
\citep{douglas2004neuronal}, and that cortical recurrent networks
provide the substrate for temporal integration and working memory
\citep{wang2001synaptic}.

\subsection{Layer-Free and Modular Architectures Are Not Equivalent}
\label{sec:disc-modular-comparison}
% ------------------------------------------------------------

In earlier work, we studied a modular variant of \modnet{}, where nodes
were grouped into a fixed number of modules. The present work is fully
layer-free: there are no modules. This was a deliberate departure, and
the results validate it.

Table~\ref{tab:modular-flat} compares the two variants on a subset of
tasks for which both results are available.

\begin{table}[h]
  \centering
  \small
  \caption{Comparison of layer-free and modular variants. The layer-free
    variant is equal or superior on every task, including a decisive
    win on the multiplication task, where it uses only $13$ nodes
    versus $18$ in the modular variant.}
  \label{tab:modular-flat}
  \begin{tabular}{llll}
    \toprule
    \textbf{Task} & \textbf{Modular} & \textbf{Layer-free} & \textbf{Winner} \\
    \midrule
    XOR-2          & 100.0\%          & 100.0\%          & tie \\
    Parity-4       & 100.0\%          & 100.0\%          & tie \\
    Parity-6       & 93.75\%          & \textbf{95.3\%}  & layer-free \\
    $\sin(2\pi x)$ & \textbf{RMSE $= 0.0439$} & RMSE $= 0.0465$ & modular \\
    $x^2 + y^2$    & RMSE $= 0.0215$  & \textbf{RMSE $= 0.0143$} & layer-free \\
    Multiplication & RMSE $= 0.0240$  & \textbf{RMSE $= 0.0254$ (13 nodes)} & layer-free \\
    \bottomrule
  \end{tabular}
\end{table}

Two observations are worth noting.

First, the layer-free variant reaches the same \emph{or better} accuracy
on every task, including the difficult Parity-6 problem where the
modular variant was stuck at $93.75\%$.

Second, the layer-free variant reaches solutions with \emph{fewer}
nodes on the continuous tasks. Multiplication, in particular, is solved
with only $13$ nodes in the layer-free variant, versus $18$ in the
modular variant.

We interpret this as evidence that modular scaffolds, far from helping,
can \emph{restrict} the search when the task does not require modular
structure. The evolution must work within the modular constraint, and
when the task is better served by a non-modular solution, the constraint
becomes a liability.

This is consistent with the ``no free lunch'' principle
\citep{wolpert1997no}, applied at the level of architectural priors: any
inductive bias that helps on some tasks will hurt on others. Modularity
is such a bias.

\subsection{Sparsity Is the Norm}
\label{sec:disc-sparsity}
% ------------------------------------------------------------

Evolved networks use a small fraction of the available connections.
Density, defined as $E / N^2$ where $E$ is the number of active
connections and $N$ is the number of nodes, ranges from $27\%$ to
$67\%$ depending on the task. This sparsity was not imposed by any
regularizer; it emerged from the pruning mechanism, which removes nodes
whose activity is either too low ($< \alpha_{\text{lo}}$) or too high
($> \alpha_{\text{hi}}$), and nodes that lie on no input-to-output path.

The observed sparsity is consistent with biological neural networks,
where the vast majority of potential synapses are absent
\citep{chklovskii2004wiring, cherniak1994component}. It is also
consistent with the finding of \citet{landassuri2012evolution} that
pure-modular architectures emerge at low connectivity values: sparse
connectivity and modular structure are related structural properties.

\subsection{Growth Adapts to Task Complexity}
\label{sec:disc-growth}
% ------------------------------------------------------------

The initial network has $26$ nodes. Over evolution, networks either
shrank, grew, or stayed approximately the same:

\begin{itemize}
  \item \textbf{Shrinking:} XOR-2 ($26 \to 8$), XOR-3 ($26 \to 9$),
    Parity-4 ($26 \to 10$), Multiplication ($26 \to 13$).
  \item \textbf{Growing:} $\sin(2\pi x)$ ($26 \to 18$), Gaussian
    ($26 \to 27$).
  \item \textbf{Nearly stable:} Parity-6, Majority-5, Full Adder,
    2-bit Adder, Comparator, MUX, Sort-4.
\end{itemize}

Boolean tasks tend to converge to small networks, while continuous
tasks require larger networks. The pattern is consistent with the
intuition that Boolean functions have low algorithmic complexity, while
continuous functions require more expressive approximations.

An interesting sub-pattern: the digital-circuit tasks (full adder,
2-bit adder, comparator, MUX, sort-4) do \emph{not} shrink
significantly. They remain near $26$ nodes, even after substantial
pruning. We interpret this as evidence that these tasks require a
minimum \emph{width} rather than depth: the network needs a certain
number of parallel nodes to resolve the input combinations, but does
not need many layers of abstraction.

This is consistent with the known circuit complexity of these tasks.
A $4$-bit sorting network requires $\Omega(n \log n)$ comparators in the
worst case \citep{knuth1998sorting}, which corresponds to a nontrivial
number of parallel nodes. Our evolved networks appear to respect this
bound.

% ------------------------------------------------------------
\section{Relationship to Related Work}
\label{sec:disc-related}
% ------------------------------------------------------------

\subsection{Comparison with NEAT and HyperNEAT}
\label{sec:disc-neat}
% ------------------------------------------------------------

NEAT \citep{stanley2002evolving} and HyperNEAT
\citep{stanley2009hypercube} are the most closely related works to
\modnet{}. All three use evolutionary search to discover network
topology, and all three have been shown to be effective on a wide range
of tasks. However, there are four important differences.

First, NEAT and HyperNEAT both use layered networks. NEAT begins with a
network of input and output nodes connected by a small number of links,
and grows by adding nodes that form a new layer between existing nodes.
HyperNEAT operates on a substrate that is typically a grid of neurons
with predefined geometric structure. \modnet{}, by contrast, has no
layers and no predefined substrate: any node may connect to any other,
and the network's geometry is not constrained.

Second, NEAT and HyperNEAT do not have a pruning mechanism. They grow
monotonically. Our results show that aggressive pruning is essential:
it removes redundancy that would otherwise accumulate, and it allows
the network to shrink when the task does not require a large size.
This is consistent with the biological observation that neural
development involves both growth and regressive events such as
synaptic pruning \citep{chklovskii2004wiring}.

Third, \modnet{} has no modular structure. This is a departure from
even our own earlier modular variant, and it distinguishes
\modnet{} from Cellular Encoding \citep{gruau1994neural}, LNDP
\citep{plantec2024evolving}, and Neural CA \citep{mordvintsev2020growing},
all of which have some implicit or explicit modular organization.

Fourth, \modnet{} is evaluated on a benchmark suite of sixteen tasks
spanning Boolean logic, digital arithmetic, sorting networks, and
continuous regression. To our knowledge, no prior layer-free
evolutionary architecture has been demonstrated on such a broad range
of tasks.

\subsection{Comparison with Self-Organizing Systems}
\label{sec:disc-self-org}
% ------------------------------------------------------------

The LNDP of \citet{plantec2024evolving} and the GRN-SO-RC of
\citet{hajizada2024developmental} are the closest self-organizing
systems to \modnet{}. Both grow networks from a single cell, and both
use local rules to control the growth process. However, both rely on
gradient-based learning for the weights, whereas \modnet{} is fully
gradient-free. This is a significant difference: gradient-based
learning requires a differentiable activation function, which imposes
constraints on the types of neurons that can be used. \modnet{}'s use
of the Heaviside step function is only possible because the weights are
evolved rather than trained.

The GRN-SO-RC also uses a modular structure, but its modules are
predetermined by the reservoir computing framework. In \modnet{}, there
are no modules at all, and the entire structure emerges from the
evolutionary process.

\subsection{Comparison with Differentiable Sorting}
\label{sec:disc-diff-sort}
% ------------------------------------------------------------

Differentiable sorting networks \citep{petersen2022deep, grover2019differentiable}
solve the sorting problem by relaxing the discrete comparator into a
continuous operator. This makes the network trainable by gradient
descent, but it requires a predefined architectural template---typically
a bitonic or odd-even sorting network---whose size grows as
$O(n \log^2 n)$.

\modnet{}, by contrast, evolves a sorting network of size $26$ for
$n = 4$ inputs. The theoretical minimum for a sorting network on $4$
inputs is $5$ comparators (Knuth 1998), which corresponds to a
network of at most $5$ nodes in a fully optimized design. Our evolved
network uses more nodes but discovers the sorting function without any
template and without gradient information. This is a different
trade-off: fewer architectural priors, more nodes.

We view this as a complementary approach: differentiable sorting is
superior when a template is available and gradient information is
useful; \modnet{} is preferable when the template is unknown and the
task must be solved from scratch.

% ------------------------------------------------------------
\section{Limitations}
\label{sec:disc-limitations}
% ------------------------------------------------------------

\subsection{Limited Scale}
\label{sec:disc-scale}
% ------------------------------------------------------------

The largest network in our experiments has $27$ nodes. This is small by
modern standards: even a modest feedforward network for a real-world
task typically has thousands of parameters. The primary reason for this
limitation is computational: the fitness evaluation is $O(E \cdot T)$
per network, and the evolutionary algorithm performs $O(P \cdot G)$
evaluations, so the total cost grows with the size of the network. On a
phone CPU, networks with more than a few hundred nodes become
impractical.

This is not a fundamental limitation of the architecture, but a
practical one. With more computational resources (a GPU cluster, a
distributed search), \modnet{} could be scaled to much larger networks.
However, scaling brings new challenges: the search space grows
exponentially with the number of nodes, and the evolutionary algorithm
may struggle to find good solutions in higher dimensions. Addressing
these challenges is a direction for future work.

\subsection{Limited Task Diversity}
\label{sec:disc-tasks}
% ------------------------------------------------------------

The sixteen tasks we considered all involve low-dimensional inputs
($2$--$8$ bits) and outputs ($1$--$8$ bits). This is far from the
dimensionality of real-world problems, which often involve images,
text, or audio with thousands or millions of dimensions. We have not
demonstrated that \modnet{} scales to these regimes, and there are
reasons to be cautious.

A more fundamental concern is the encoding. Our current encoding
requires quantizing continuous values to a fixed number of bits. For a
$1000$-dimensional input with $8$-bit precision per dimension, the
input would be $8000$ bits, and even a small network would need
thousands of nodes to process them. Whether \modnet{} can handle this
scale is an open question.

\subsection{Two Failure Cases}
\label{sec:disc-failures}
% ------------------------------------------------------------

Two tasks in our suite did not reach their success threshold:
Parity-6 ($95.3\%$) and $\sin(4\pi x)$ (RMSE $= 0.0789$).

For Parity-6, the network reached a local optimum with $8$ nodes and
five self-loops, and could not escape it despite additional growth and
pruning. This is a real failure, not a technical artifact. We discuss
several hypotheses for why it occurred in
Section~\ref{sec:disc-parity6-hypotheses}.

For $\sin(4\pi x)$, the network reached an RMSE of $0.0789$, which is
above our success threshold of $0.05$ but below the alternative
threshold of $0.10$ that is sometimes used for high-frequency
regression tasks. We consider this a partial success rather than a
failure.

\subsection{No Theoretical Guarantees}
\label{sec:disc-theory}
% ------------------------------------------------------------

We have not proven any theoretical properties of \modnet{}. In
particular, we have not shown that the evolutionary algorithm converges
to a global optimum, nor that the architecture is universal in the
sense of Turing completeness. The empirical results are encouraging,
but they do not constitute a proof.

This is a common situation in neuroevolution: most algorithms in this
family are evaluated empirically rather than theoretically. However,
the lack of guarantees means that \modnet{}'s behavior in new tasks is
unpredictable: it may succeed brilliantly, or it may fail
catastrophically, and we do not have a principled way to predict which
will happen.

\subsection{Single Seed}
\label{sec:disc-single-seed}
% ------------------------------------------------------------

Our reported results use a single seed per task, with the seed reset to
a common value at the start of each task. This ensures reproducibility
across runs, but it does not provide statistical error bars. A
multi-seed study---running each task with five to ten different seeds
and reporting mean and standard deviation---would strengthen the
empirical claims.

We leave multi-seed analysis to future work, but we note that the
observed patterns (recurrence scaling, sparsity, growth regimes) are
consistent across all sixteen tasks, which suggests they are robust
patterns rather than seed-specific artifacts.

% ------------------------------------------------------------
\section{Threats to Validity}
\label{sec:disc-threats}
% ------------------------------------------------------------

\subsection{Internal Validity}
\label{sec:disc-internal}
% ------------------------------------------------------------

The main internal threat to validity is the choice of fitness function.
We used mean squared error for all tasks, which is a standard but not
universal choice. For Boolean tasks, a different fitness function (such
as classification accuracy) might lead to different results, because
the gradients of the two functions differ. For regression tasks, a
robust loss (such as Huber loss) might outperform MSE in the presence
of outliers.

A second internal threat is the choice of activation function. We used
the Heaviside step function, which is the simplest threshold activation.
Other choices (such as sigmoid or ReLU) are possible but would require
a different learning algorithm (since the step function is not
differentiable). It is not clear whether \modnet{} would perform
similarly with a different activation.

\subsection{External Validity}
\label{sec:disc-external}
% ------------------------------------------------------------

The main external threat to validity is the choice of tasks. Our
sixteen tasks are all small and synthetic. They were chosen to span a
range of difficulties, from XOR to 4-bit sorting network and from a
single sinusoid to multiplication. However, they are far from the tasks
that are typically used to evaluate neural networks. In particular,
none of our tasks involves structured inputs (images, sequences,
graphs), and none requires long-range dependencies.

It is possible that \modnet{} would perform well on some real-world
tasks and poorly on others, and we do not have a principled way to
predict which. This is a general limitation of empirical evaluations
in machine learning, and it is not specific to \modnet{}.

\subsection{Construct Validity}
\label{sec:disc-construct}
% ------------------------------------------------------------

The main construct threat to validity is the interpretation of the
structural metrics we use. Our measure of recurrence (the number of
self-loops) captures only one type of recurrence. Longer cycles
($i_1 \to i_2 \to \cdots \to i_\ell \to i_1$) are not counted, and
these may be just as important for the network's function. A more
thorough analysis would count all cycles, not just self-loops.

Similarly, our measure of sparsity (the ratio of active to total
connections) treats all connections as equally important. In a
biological network, some connections are strong and some are weak, and
the strength matters. Our metric cannot distinguish between a network
with many weak connections and a network with a few strong ones.

% ------------------------------------------------------------
\section{Hypotheses for the Parity-6 Failure}
\label{sec:disc-parity6-hypotheses}
% ------------------------------------------------------------

Parity-6 is the only Boolean task in our suite that did not reach exact
accuracy. We consider several hypotheses for why it failed.

\subsection{Hypothesis 1: Insufficient Time Steps}
\label{sec:disc-h1}
% ------------------------------------------------------------

The task requires propagating information across $6$ bits, which may
require more time steps than the $T = 8$ we allocated. Parity-4 was
solved with $T = 6$, so Parity-6 may require $T \geq 10$.

This hypothesis predicts that increasing $T$ to $10$ or $12$ should
allow the network to solve the task. We consider this the most likely
hypothesis.

\subsection{Hypothesis 2: Local Optimum with Under-Parameterization}
\label{sec:disc-h2}
% ------------------------------------------------------------

The best Parity-6 network reached $95.3\%$ accuracy with only $8$ nodes,
while the successful Parity-4 network used $10$ nodes. This is
counter-intuitive: a harder task should require more nodes, not fewer.
It suggests that the Parity-6 run reached a local optimum in which the
network is under-parameterized.

This hypothesis predicts that forcibly increasing the minimum network
size for Parity-6 should help. A possible fix is to disable pruning for
a few hundred generations, allowing the network to grow before any
shrinking is attempted.

\subsection{Hypothesis 3: Insufficient Population Diversity}
\label{sec:disc-h3}
% ------------------------------------------------------------

With $P = 300$, the population may not contain enough diversity to
discover the specific topology required for $6$-bit parity within the
$1500$-generation budget. Parity-6 is known to be exponentially harder
than Parity-4 for bounded-depth threshold circuits
\citep{impagliazzo1997size}, so the search space is correspondingly
larger.

This hypothesis predicts that increasing $P$ to $500$ or $1000$ should
help. A preliminary test with $P = 500$ reached $95.3\%$ accuracy
(versus $93.75\%$ with $P = 300$), so this hypothesis is partially
supported.

% ------------------------------------------------------------
\section{Future Work}
\label{sec:disc-future}
% ------------------------------------------------------------

\subsection{Scaling to Larger Networks}
\label{sec:disc-future-scale}
% ------------------------------------------------------------

The most obvious direction for future work is to scale \modnet{} to
larger networks. There are several ways to approach this.

First, we can parallelize the fitness evaluation. Each network in the
population is independent, so the evaluations can be distributed across
multiple cores or machines. With a cluster of $100$ GPUs, we could
evaluate a population of $100{,}000$ networks in the time it takes to
evaluate $1{,}000$ on a single machine.

Second, we can use a more efficient encoding. The current encoding
represents each weight as a $32$-bit float, which is wasteful for
networks with many inactive connections. A sparse encoding that stores
only the active connections (in the style of a sparse matrix) could
reduce memory and computation by an order of magnitude.

Third, we can use a more sophisticated search algorithm. The
evolutionary algorithm we use is simple; more advanced methods such as
CMA-ES \citep{hansen2001completely}, differential evolution
\citep{storn1997differential}, or Bayesian optimization
\citep{snoek2012practical} might converge faster. The choice of
algorithm is orthogonal to the architecture: \modnet{} can be combined
with any evolutionary search method.

\subsection{Multi-Seed Analysis}
\label{sec:disc-future-seed}
% ------------------------------------------------------------

A multi-seed study is a natural extension of the present work. Running
each task with five to ten seeds and reporting mean and standard
deviation would strengthen the empirical claims. It would also allow
us to test whether the observed patterns (recurrence scaling, sparsity,
growth regimes) hold across seeds.

\subsection{Extending to Structured Inputs}
\label{sec:disc-future-structured}
% ------------------------------------------------------------

Our current encoding assumes that inputs are unstructured vectors of
bits. Many real-world tasks involve structured inputs: images have
spatial structure, sequences have temporal structure, and graphs have
relational structure. Extending \modnet{} to handle these inputs would
require incorporating structural priors into the architecture.

One approach is to use a convolutional or graph-based preprocessing
layer to extract local features, and then feed these features to a
\modnet{} network. This is analogous to the way that modern vision
models use a convolutional backbone followed by a fully connected head.
The difference is that the head in our case is a layer-free recurrent
network rather than a stack of dense layers.

Another approach is to make \modnet{} itself structure-aware, by
allowing nodes to have ``receptive fields'' that are defined by the
input structure. This is closer to the approach taken by HyperNEAT
\citep{stanley2009hypercube}, where the geometry of the substrate
encodes the structure of the input. Extending this idea to the
layer-free setting is an interesting direction.

\subsection{Combining Evolution with Gradient Descent}
\label{sec:disc-future-hybrid}
% ------------------------------------------------------------

A natural extension of \modnet{} is to combine evolutionary topology
search with gradient-based weight optimization. In this hybrid
approach, the evolutionary algorithm would discover the topology, and
gradient descent would fine-tune the weights. This is analogous to the
approach taken by \citet{real2019regularized} for image classification,
but applied to the layer-free setting.

The challenge is that the Heaviside activation is not differentiable,
so standard backpropagation cannot be applied directly. One solution is
to use a differentiable surrogate for the activation during training,
and then use the original Heaviside activation during evaluation. This
is the approach taken by spiking neural networks
\citep{neftci2019surrogate}. Another approach is to replace the
Heaviside with a sigmoid during a ``polishing'' phase, and then revert
to the Heaviside once the weights have converged.

We have explored the second approach in preliminary experiments, but
the results were not encouraging: the gradient information was not
informative enough to improve the weights significantly. A more
thorough investigation is needed to determine whether this approach can
work in practice.

\subsection{Theoretical Analysis}
\label{sec:disc-future-theory}
% ------------------------------------------------------------

Finally, a more thorough theoretical analysis of \modnet{} would be
valuable. There are several questions we would like to answer:

\begin{itemize}
  \item What is the computational power of \modnet{}? Is it Turing
    complete? Under what conditions on the number of nodes and time
    steps?
  \item What is the relationship between the evolved structure and the
    task? Can we predict the number of nodes required for a given task?
    Can we characterize the class of tasks that \modnet{} can solve
    efficiently?
  \item What is the convergence rate of the evolutionary algorithm?
    How many generations are required to reach a given level of
    accuracy? How does this depend on the task?
  \item Is the recurrence count a fundamental invariant? Are there
    tasks for which the minimum recurrence is provably $\Omega(f(n))$
    for some function $f$?
\end{itemize}

Answering these questions would require a combination of empirical
studies and theoretical analysis, and would likely require tools from
several fields (evolutionary computation, complexity theory, and
statistical learning theory).

% ------------------------------------------------------------
\section{Broader Implications}
\label{sec:disc-broader}
% ------------------------------------------------------------

Beyond the technical results, this work has broader implications for
the way we think about neural architecture design.

The dominant paradigm in modern machine learning is to design
architectures by hand, guided by intuition and empirical tuning. This
has been extremely successful, but it has also led to a proliferation
of architectures that are highly specialized for particular tasks. The
transformer, for example, is optimized for sequences; the convolutional
network is optimized for images; the recurrent network is optimized
for time series. Each architecture comes with its own set of
hyperparameters, its own training procedure, and its own failure modes.

\modnet{} suggests an alternative: a single, general-purpose
architecture that can be adapted to any task through evolution. The
same node model, the same fitness function, and the same evolutionary
algorithm solved Boolean logic, digital arithmetic, sorting networks,
and continuous regression with no task-specific modifications.

This is not to say that \modnet{} is better than specialized
architectures on their home turf---it is not, at least not at the scale
we have tested. But it does suggest that architecture design can be
automated, and that the resulting architectures need not be
specialized.

A second implication concerns the relationship between architectural
priors and task difficulty. Our comparison of layer-free and modular
variants suggests that priors that help on some tasks hurt on others.
This is a specific instance of the ``no free lunch'' theorem
\citep{wolpert1997no}: no single architecture is best for all problems,
but a single architecture combined with a powerful search algorithm
might be able to solve a wide range of problems by discovering
task-specific structures on the fly. This is the essence of
meta-learning, and \modnet{} is a step in that direction.

% ------------------------------------------------------------
\section{Summary}
\label{sec:disc-summary}
% ------------------------------------------------------------

This chapter has interpreted the results of Chapter~\ref{ch:results}
and discussed their limitations. The main conclusions are:

\begin{enumerate}
  \item \textbf{Layer-free computation is feasible for digital logic.}
    Eleven of eleven Boolean and digital-circuit tasks were solved to
    exact $100\%$ accuracy, including full adder, 2-bit adder,
    comparator, MUX, and 4-bit sorting network.

  \item \textbf{Recurrence is a structural necessity.} The number of
    self-loops scales with the temporal complexity of the task,
    ranging from one (XOR-2) to fourteen (Gaussian).

  \item \textbf{Layer-free is equal or superior to modular.} The
    layer-free variant is never worse than the modular variant on any
    task, and is decisively better on the multiplication task.

  \item \textbf{The architecture has limits.} Parity-6 was not solved
    exactly within our computational budget, and there is no reason to
    expect \modnet{} to scale indefinitely.

  \item \textbf{There is much room for improvement.} Multi-seed
    analysis, structured inputs, hybrid evolution-gradient methods,
    and theoretical analysis are all promising directions for future
    work.
\end{enumerate}

These observations motivate the concluding remarks in
Chapter~\ref{ch:conclusion}.

% ============================================================
%  Bibliography for this chapter
% ============================================================
